% Generated draft: requires mathematical and source review.
\documentclass[11pt]{article}
\usepackage[margin=25mm]{geometry}
\usepackage{fontspec}
\setmainfont{Latin Modern Roman}
\usepackage{amsmath,amssymb,mathtools,unicode-math}
\setmathfont{Latin Modern Math}
\usepackage{enumitem,needspace,xurl,hyperref}
\setlength{\parindent}{0pt}
\setlength{\parskip}{5pt}
\setlength{\emergencystretch}{3em}
\newcommand{\N}{\mathbb{N}}
\newcommand{\Z}{\mathbb{Z}}
\newcommand{\Q}{\mathbb{Q}}
\newcommand{\R}{\mathbb{R}}
\newcommand{\C}{\mathbb{C}}
\newcommand{\F}{\mathbb{F}}
\newcommand{\A}{\mathbb{A}}
\newcommand{\PP}{\mathbb{P}}
\newcommand{\E}{\mathbb{E}}
\newcommand{\eps}{\varepsilon}
\newcommand{\dd}{\,\mathrm{d}}
\begin{document}
\section*{UnsolvedMath ID: 800793}
\textbf{Source problem code:} catalogue entry
\subsection{Definitions and mathematical statement}
Let \(n\) be a positive integer, and let \(f\colon(\C^n,0)\to(\C^n,0)\) be a holomorphic germ tangent to the identity, meaning \(f(0)=0\) and \(Df(0)=I_n\). For a nonidentity such germ, write
\[
f(z)=z+P_{\nu}(z)+\sum_{j>\nu}P_j(z),\qquad \nu\geq2,
\]
where each \(P_j\colon\C^n\to\C^n\) is homogeneous of degree \(j\) and \(P_{\nu}\not\equiv0\). A characteristic direction is a projective class \([v]\in\PP^{n-1}(\C)\), with \(v\neq0\), for which
\[
P_{\nu}(v)=\lambda v
\]
for some \(\lambda\in\C\). Such a direction is called degenerate if \(P_{\nu}(v)=0\). Thus the condition that all characteristic directions are degenerate is
\[
\forall v\in\C^n\setminus\{0\},\qquad
P_{\nu}(v)\in\C v\ \Longrightarrow\ P_{\nu}(v)=0.
\]
A formal change of coordinates is a formal power series map \(H\in\C[[z_1,\ldots,z_n]]^n\) with \(H(0)=0\) and invertible linear part. Two germs are formally conjugate under such a change when
\[
H\circ f=g\circ H.
\]
A formal invariant is a quantity constant under the chosen formal conjugacy relation; a formal normal form is a specified representative of a formal conjugacy class.

\textbf{Problem (Abate, OP15, 2011).} Find formal normal forms and formal invariants for holomorphic germs tangent to the identity having only degenerate characteristic directions.

The supplied formulation does not specify the allowed dimensions, whether the identity germ is included, or whether coordinate changes may have arbitrary invertible linear part or must themselves be tangent to the identity. The definitions above explain the nonidentity case in finite complex dimension and the usual unrestricted formal conjugacy relation; these unresolved scope and equivalence conventions require confirmation before a definitive classification target can be fixed. The source reports partial results under the reference labels [BM], [AT2], and [AR], without supplying their bibliographic entries.
\subsection{Short English statement}
Abate's OP15 asks for formal normal forms and invariants of holomorphic germs tangent to the identity whose characteristic directions are all degenerate. The dimension, treatment of the identity germ, and permitted formal coordinate changes need clarification. The supplied registry labels the problem open and reports partial results in [BM], [AT2], and [AR]; these assessments have not been independently verified.
\subsection{Sources}
\begin{itemize}
\item Local UnsolvedMath registry, permanent ID \(800793\), exact source problem code \texttt{catalogue entry}, title: \emph{Local Discrete Holomorphic Dynamics --- Open Problem OP15}. The supplied record attributes the problem to Abate, with proposed year 2011, and labels its status open.
\item Abate, \emph{Open Problems in Local Discrete Holomorphic Dynamics} (2011), Problem OP15. Supplied source URL: \url{https://arxiv.org/abs/1107.4876}. The supplied provenance records access on 2026-10-16 and extraction by \texttt{arxiv\_latex}; this access information is supplied provenance, not a claim of independent source consultation.
\item The supplied statement cites partial results as [BM], [AT2], and [AR]. Their bibliographic entries and individual links were not supplied, so these references remain unresolved.
\end{itemize}
\end{document}
